\section{Introduction}\label{sec:intro}

Conformal prediction turns the output of any trained predictor into a prediction set that contains the true label with a chosen probability, and it needs no assumption on the data distribution except exchangeability \cite{vovk2005algorithmic,shafer2008tutorial,angelopoulos2023conformal}. In classification, the set is built from a score that measures how badly a candidate label agrees with an input. The coverage guarantee holds for every choice of score, but the choice decides the size of the sets: a poorly chosen score gives sets that are valid but too large to be useful.

Often several predictors score each candidate class, for example classifiers with different models or different input features. The scores of one candidate then form a vector, and a single score has to be obtained from it. A fixed rule such as the mean treats all predictors in the same way, although they can differ in reliability and their scores can be dependent. Methods for conformal aggregation mostly combine the prediction sets of the individual predictors \cite{gasparin2024merging,yang2024selection}, or project the score vector on one weighted direction \cite{luo2024weighted}. Ochoa Rivera et al.~\cite{ochoarivera2024conformal} instead work in score space. They define a quantile envelope for the multivariate score as the intersection of half-spaces along random directions, learn its shape on one part of the calibration data and scale it on the other part, which gives marginal coverage.

In this paper we build on this idea in a classification setting. The envelope is learned separately for each class, from the score vectors of the samples of that class, so that the coverage statement holds for each class; this is known as Mondrian conformal prediction \cite{vovk2005algorithmic}. We present three constructions that share the same two stage calibration: a collapsed envelope, which averages the scores, a radial envelope, in which the boundary changes with the direction of the score vector, and a strip envelope, which bounds each coordinate given the value of another coordinate. The radial and strip envelopes are not intersections of half-spaces, and all three constructions use only the observed scores of a sample, so a score vector with missing entries can still be assigned a prediction set.

Our contributions are the three constructions together with a class-conditional coverage statement and its assumptions (Sections~\ref{sec:setup} and~\ref{sec3:constructions}), an open-source Python package that implements them (Section~\ref{sec:software}), and experiments on simulated data with missing scores and on football player data (Section~\ref{sec:experiments}). In these experiments the coverage is close to the nominal level for all three constructions, and we do not claim that one of them is better than the others in general. The last section discusses the limitations of our work.